%======================================================================
% CAPÍTULO 14 - MOTOR DE CÁLCULO DE NÓMINA
%======================================================================
\chapter{Motor de Cálculo de Nómina}

\bigskip

El motor de cálculo constituye el corazón de Sentinel. Este capítulo detalla la arquitectura del motor, las fórmulas implementadas y la secuencia de cálculos.

\bigskip

\section{Arquitectura del Motor}

El motor de cálculo se divide en dos componentes fundamentales:

\bigskip

\begin{enumerate}
    \item Motor Rhai: Evalúa fórmulas de primas definidas en la base de datos.
    \item Calculadora Nativa: Calcula componentes de nómina en Rust nativo.
\end{enumerate}

\bigskip

El motor de cálculo implementa un enfoque de pipeline secuencial, donde cada paso depende del resultado del anterior, garantizando la coherencia de los cálculos.

\bigskip

\section{Secuencia de Cálculos}

El cálculo de nómina se ejecuta en el siguiente orden para garantizar la correcta dependencia entre valores:

\bigskip

\begin{enumerate}
    \item Tiempo de Servicio
    \item Antigüedad Grado
    \item Sueldo Base
    \item Calcular Primas (Rhai)
    \item Sueldo Mensual
    \item Alicuota Aguinaldo
    \item Alicuota Vacaciones
    \item Sueldo Integral
    \item Asignación Antigüedad
    \item Garantías
    \item Días Adicionales
    \item No Depositado Banco
\end{enumerate}

\bigskip

\section{Fórmulas de Cálculo}

\subsection{Cálculo de Antigüedad}

El sistema calcula el tiempo de servicio utilizando aritmética de fechas con precisión de días:

\[
\text{Antigüedad} = \text{Fecha\_Actual} - \text{Fecha\_Ingreso}
\]

\[
\text{Antigüedad\_Grado} = \text{Fecha\_Actual} - \text{Fecha\_Último\_Ascenso}
\]

\bigskip

\subsection{Sueldo Mensual}

\[
\text{Sueldo Mensual} = \text{Sueldo Base} + \sum \text{Primas}
\]

Donde:
\begin{itemize}
    \item Sueldo Base: Obtenido de la tabla de sueldos (Directiva).
    \item Primas: Suma de todas las primas calculadas por el Motor Rhai.
\end{itemize}

\bigskip

\subsection{Alicuota de Aguinaldo}

\[
\text{Alicuota Aguinaldo} = \frac{\text{Días\_Aguinaldo} \times \text{Sueldo Mensual}}{30 \times 12}
\]

\bigskip

\begin{center}
\begin{tabular}{ll}
\toprule
\textbf{Condición} & \textbf{Días de Aguinaldo} \\
\midrule
Retiro $<$ 2016 & 90 \\
2016-10-01 a 2016-12-31 & 105 \\
Retiro $\ge$ 2017-01-01 & 120 \\
\bottomrule
\end{tabular}
\end{center}

\bigskip

\subsection{Alicuota de Vacaciones}

\[
\text{Alicuota Vacaciones} = \frac{\text{Días\_Vacaciones} \times \text{Sueldo Mensual}}{30 \times 12}
\]

\bigskip

\begin{center}
\begin{tabular}{ll}
\toprule
\textbf{Tiempo de Servicio} & \textbf{Días de Vacaciones} \\
\midrule
1-14 años & 40 \\
15-24 años & 45 \\
25+ años & 50 \\
\bottomrule
\end{tabular}
\end{center}

\bigskip

\subsection{Sueldo Integral}

\[
\text{Sueldo Integral} = \text{Sueldo Mensual} + \text{Alicuota Vacaciones} + \text{Alicuota Aguinaldo}
\]

\bigskip

\subsection{Asignación de Antigüedad}

\[
\text{Asignación Antigüedad} = \text{Sueldo Integral} \times \text{Tiempo de Servicio (años)}
\]

\bigskip

\subsection{Garantías}

\[
\text{Garantías} = \frac{\text{Sueldo Integral}}{30} \times 15
\]

\bigskip

\subsection{Días Adicionales}

\[
\text{Factor} = \min(\text{Tiempo de Servicio}, 15)
\]

\[
\text{Días Adicionales} = \frac{\text{Sueldo Mensual}}{30} \times 2 \times \text{Factor}
\]

\bigskip

\subsection{No Depositado en Banco}

\[
\text{No Depositado} = \text{Asignación Antigüedad} - \text{Depósito Banco} - \text{Garantías} - \text{Días Adicionales}
\]

\bigskip

\section{Motor Rhai: Evaluación de Primas}

Las primas se definen como expresiones Rhai en la base de datos, permitiendo modificación sin recompilación.

\bigskip

Ejemplo de fórmula Rhai para prima de antigüedad:
\begin{verbatim}
prima_antiguedad = if antiguedad > 0 { 
    (sueldo_base * antiguedad) / 100.0 
} else { 0.0 }
\end{verbatim}

\bigskip

Ejemplo de fórmula para prima de hijos:
\begin{verbatim}
prima_hijos = if n_hijos > 0 {
    (sueldo_base * n_hijos * 0.05)
} else { 0.0 }
\end{verbatim}

\vfill
\newpage
